% Generated draft: requires mathematical and source review.
\documentclass[11pt]{article}
\usepackage[margin=25mm]{geometry}
\usepackage{fontspec}
\setmainfont{Latin Modern Roman}
\usepackage{amsmath,amssymb,mathtools,unicode-math}
\setmathfont{Latin Modern Math}
\usepackage{enumitem,needspace,xurl,hyperref}
\setlength{\parindent}{0pt}
\setlength{\parskip}{5pt}
\setlength{\emergencystretch}{3em}
\newcommand{\N}{\mathbb{N}}
\newcommand{\Z}{\mathbb{Z}}
\newcommand{\Q}{\mathbb{Q}}
\newcommand{\R}{\mathbb{R}}
\newcommand{\C}{\mathbb{C}}
\newcommand{\F}{\mathbb{F}}
\newcommand{\A}{\mathbb{A}}
\newcommand{\PP}{\mathbb{P}}
\newcommand{\E}{\mathbb{E}}
\newcommand{\eps}{\varepsilon}
\newcommand{\dd}{\,\mathrm{d}}
\begin{document}
\section*{UnsolvedMath ID: 20001296}
\textbf{Source problem code:} catalogue entry
\subsection{Definitions and mathematical statement}
The precise finite-chain formulation supplied by the research summary is as follows. Let $P$ be the transition matrix of a finite irreducible discrete-time Markov chain on $\Omega$, with stationary probability $\mu$, reversible in the sense that $\mu(x)P(x,y)=\mu(y)P(y,x)$, and lazy in the sense that $P(x,x)\geq1/2$ for every $x$. Define
\[
 \|\nu-\mu\|_{\mathrm{TV}}=\tfrac12\sum_{y\in\Omega}|\nu(y)-\mu(y)|,
 \qquad
 \tau_{\mathrm{TV}}(\delta)=\min\{t\in\Z_{\geq0}:\max_x\|P^t(x,\cdot)-\mu\|_{\mathrm{TV}}\leq\delta\}.
\]
For $A\subseteq\Omega$ put $T_A=\inf\{t\geq0:X_t\in A\}$, and, for $0<\beta\leq1$, put
\[
 T(\beta)=\max_{x\in\Omega,\ A\subseteq\Omega:\mu(A)\geq\beta}\E_xT_A.
\]
Yuval Peres' question asks for a universal constant $c>0$, independent of the chain and state-space size, such that
\[
 \tau_{\mathrm{TV}}(1/2)\leq c\,T(1/2).
\]
Lazy random walks on finite connected undirected graphs are included. The original excerpt does not explicitly specify finiteness or irreducibility, and its mixing definition leaves the worst-case starting-state quantifier implicit; these are made explicit here according to the supplied research formulation, rather than asserting an extension to arbitrary chains.

The supplied research summary labels the finite lazy reversible question solved in the literature: Griffiths, Kang, Oliveira and Patel are credited with universal comparability of $T(1/2)$ and $\tau_{\mathrm{TV}}(1/4)$, which implies the requested inequality. With
\[
 t_{\mathrm{prod}}=\max_{x\in\Omega,\ A\subseteq\Omega:\mu(A)>0}\mu(A)\E_xT_A,
\]
it reports their sharp endpoint comparison
\[
 \tfrac12 T(1/2)\leq t_{\mathrm{prod}}\leq T(1/2),
\]
and attributes the mixing-time comparison with $t_{\mathrm{prod}}$ to Aldous. The original Aldous remark has damaged placement of $\mu(A)$ relative to the maximum; the coherent product parameter above is supported by the supplied summary, not by an independently checked source.

The aid further reports the sharp special-case inequality $\tau_{\mathrm{TV}}(1/2)\leq T(1/2)$ for irreducible lazy two-state chains
\[
 P=\begin{pmatrix}1-a&a\\b&1-b\end{pmatrix},\qquad 0<a,b\leq\tfrac12,
\]
with optimal universal constant $1$, approached by $b=1/2-\varepsilon$, $a=\varepsilon^2$ as $\varepsilon\downarrow0$. It also reports failure of a corresponding universal bound using target-set mass threshold $\beta>1/2$ in place of $1/2$. These special-case and threshold claims are retained as attributed research-aid statements, without adding an unprovided counterexample or proof.
\subsection{Short English statement}
The question compares worst-case mixing with the longest expected time to hit a set carrying at least half the stationary mass. The supplied research describes an affirmative published result for finite irreducible lazy reversible chains, including the half-mass endpoint, and a sharp two-state special case. The original broad wording and damaged Aldous remark are qualified explicitly.
\subsection{Sources}
\begin{itemize}
\item Local UnsolvedMath registry, permanent ID 20001296, exact source code catalogue entry; canonical \texttt{aim-convex-geometry-notes.json}, \texttt{notes[27]}, source item 9, tag ``problem'', workshop \emph{Algorithmic convex geometry}. Source attribution: Yuval Peres; supplied verification identifies \emph{Problems from the AIM Workshop on Algorithmic Convex Geometry} (2007). Supplied URL: \url{https://aimath.org/WWN/convexgeometry/convexgeometry.pdf}.
\item Original remark credits Aldous with a mixing/hitting-time comparison, but its displayed $c\mu(A)\max_{x,A}\E T_{A,x}$ is typographically ambiguous. The supplied summary clarifies the product parameter and credits Griffiths, Kang, Oliveira and Patel with the half-mass endpoint result. No paper titles, publication dates, or full bibliographic entries for either attribution are supplied; none is invented.
\item Registry status ``solved'', research classification SOLVED-IN-LITERATURE, statement status ``exact''; attempt 1, result ``full\_solution'', problem status at run ``solved'', recorded rigor ``proved'', novelty confidence ``low'', candidate type ``sharp\_special\_case''. The source calls this a machine-generated, non-peer-reviewed research aid requiring independent verification before citation. The literature field is empty. The verification truncates after announcing its intended worst-case mixing definition; no outside source was visited.
\end{itemize}
\end{document}
